%Indique ac\'a, las herramientas de LLMs/IA usadas en su proyecto. \\En caso de no utilizar ninguna explicitar:
%\textbf{Declaro no haber utilizado ninguna herramienta de LLMs/IA en el desarrollo de este trabajo.}\\

%\noindent En caso de s\'i utilizarlo, para cada herramienta/uso, indicar:
\begin{enumerate}
\item Herramienta utilizada: Gemini Pro y NotebookLM en sus versiones gratuitas para estudiante
\item Propósito de uso: Generar las bibliografias de .bib. Revisión de redacción. Generar listados de comparación entre modelos según paper. Escritura en latex de formulas.
\item Prompts relevantes: 
\begin{itemize}
\item Generame el .bib de la referencia [X]
\item Comparame el modelo de la fuente [X] con el de la fuente [Y].
\item Dado el siguiente archivo .tex, indicame los errores ortograficos y de redacción, considerando que se debe de escribir en tercera persona.
\item Retorname en una ecuación de latex la ecuación [i] del paper [X]
\end{itemize}
\item Nivel de edición: Reescritura y revisión directa en el paper en el caso de comparaciones de modelos. Copia y pega de los .bib y formulas de latex generadas. Dependiendo del caso, las revisiones de redacción fueron seleccionadas manualmente.
\item Validación: Lectura de la referencia directa (NotebookLM entrega la sección especifica de donde saca la información).
\end{enumerate}


%Herramienta: ChatGPT (GPT-5.x)
%Uso: Generación de ideas y revisión de redacción
%Prompt principal: “Explica el algoritmo SMS-EMOA con pseudocódigo formal”
%Intervención: El texto fue adaptado, reorganizado y complementado con material de clase
%Validación: Se contrastó con [fuente X]
